% SPDX-License-Identifier: CC-BY-4.0
\documentclass[11pt,letterpaper]{article}

\usepackage[T1]{fontenc}
\usepackage{amsmath,amssymb}
\usepackage{tikz}
\usepackage{pgfplots}
\pgfplotsset{compat=1.18}
\usepackage{algorithm}
\usepackage{algpseudocode}
\usepackage{booktabs}
\usepackage{enumitem}
\usepackage[hidelinks]{hyperref}

\newcommand{\N}{\mathbb{N}}
\newcommand{\Z}{\mathbb{Z}}
\newcommand{\Q}{\mathbb{Q}}
\newcommand{\R}{\mathbb{R}}
\newcommand{\C}{\mathbb{C}}
\newcommand{\Map}[3]{#1 \colon #2 \to #3}
\newcommand{\Mod}[1]{#1\text{-}\mathrm{Mod}}
\newcommand{\GroupHom}{\mathrm{Hom}_{\mathbf{Grp}}}
\newcommand{\RingHom}{\mathrm{Hom}_{\mathbf{Ring}}}

\newtheorem{theorem}{Theorem}
\newtheorem{lemma}[theorem]{Lemma}
\newtheorem{corollary}[theorem]{Corollary}
\newtheorem{definition}[theorem]{Definition}
\newtheorem{remark}[theorem]{Remark}
\newenvironment{proof}{\noindent\textit{Proof.}}{\hfill\(\square\)}

\title{Generalized Collatz Library}
\author{}
\date{\today}

\begin{document}
\maketitle
\section*{Scalar Off-Axis Escape on Imaginary Quadratic Orders}

\begin{abstract}
Let \(R=\Z[\theta]\subset\C\), where
\(\theta^2=\tau\theta-\nu\), \(\tau,\nu\in\Z\), and \(4\nu>\tau^2\).
For any scalar integral offset \(c\), every nonreal seed of the division-first
raw map \(z\mapsto z/d\) or \(qz+c\) escapes in the complex norm if
\(|q|,|d|\ge2\) have no equal positive powers. The conclusion holds along
the full original raw sequence. The exact transverse identity, rather than
an eventual division-free tail, excludes every repeated transverse magnitude.
Integer discreteness completes the proof. Scalar coordinate matching also
gives a same-clock conjugacy between all these orders, although it is not
a ring homomorphism. A shared-prime-support example keeps dividing forever,
and a dependent pair has a primitive nonreal five-cycle.
\end{abstract}

\section{The original ring and guard}

Fix \(\tau,\nu\in\Z\) with \(4\nu>\tau^2\) and choose
\[
\theta=\frac{\tau+i\sqrt{4\nu-\tau^2}}2.
\]
Then \(\theta^2=\tau\theta-\nu\) and \(\operatorname{Im}\theta>0\).
Every element of \(R=\Z[\theta]\) has a unique expression \(x+y\theta\)
with \(x,y\in\Z\): reduction by the monic quadratic gives existence,
and equality of imaginary parts followed by real parts gives uniqueness.
For scalar parameters \(q,c,d\in\Z\), \(|d|\ge2\), define
\[
T_{q,c,d}(z)=
\begin{cases}
z/d,&z\in dR,\\
qz+c,&z\notin dR.
\end{cases}
\]
Division has precedence, including at zero; the sign of \(d\) is retained.
The Gaussian order uses \((\tau,\nu)=(0,1)\), and the Eisenstein order
uses \((\tau,\nu)=(-1,1)\). This is ring divisibility, not divisibility
in the fraction field.

% The coordinate guard and quotient are proved in the actual integral ring.
\begin{lemma}\label{Lem:Guard}
A scalar \(d\) divides \(x+y\theta\) in \(R\) exactly when \(d\mid x\)
and \(d\mid y\). Its quotient has coordinates \((x/d,y/d)\), and that
quotient is unique for \(d\ne0\). Consequently the coordinate map is
\[
(x,y)\longmapsto
\begin{cases}
(x/d,y/d),&d\mid x\text{ and }d\mid y,\\
(qx+c,qy),&\text{otherwise}.
\end{cases}
\]
\end{lemma}
\begin{proof}
If \(x+y\theta=d(u+v\theta)\), uniqueness of coordinates gives
\(x=du\) and \(y=dv\). Conversely these two integer equalities construct
the ring quotient \(u+v\theta\). If two quotients exist, subtract their
coordinate equalities and cancel the nonzero integer \(d\) in each
coordinate. No assumption about a chosen branch or a fraction-field
quotient is needed. Multiplication by a scalar and addition of a scalar
have the displayed affine coordinates.
\end{proof}

\section{An exact transverse invariant}

Write \(T^n(z)=x_n+y_n\theta\). Let \(a_n\) and \(b_n\) count affine and
division instructions among the first \(n\) raw ticks.

% This prefix formula retains every guard and the original raw clock.
\begin{lemma}\label{Lem:Product}
For every \(n\ge0\),
\[
a_n+b_n=n,\qquad d^{b_n}y_n=q^{a_n}y_0.
\]
If \(q\ne0\) and \(y_0\ne0\), then \(y_n\ne0\) for all \(n\).
\end{lemma}
\begin{proof}
At time zero both counts vanish and the product identity is \(y_0=y_0\).
An affine instruction sends \(y_n\) to \(qy_n\) and adds one to \(a_n\),
so multiplication by \(q\) preserves the identity. A division instruction
is legal only if \(d\mid y_n\); it sends \(y_n\) to its exact signed
quotient and adds one to \(b_n\). The equality
\(d(y_n/d)=y_n\) preserves the identity. These two cases are exhaustive,
regardless of how \(x_n\) affects the common guard.
The right side \(q^{a_n}y_0\) is nonzero under the stated hypotheses,
so the left side cannot have \(y_n=0\).
\end{proof}

\begin{definition}\label{Def:Independent}
Positive integers \(Q,D\ge2\) are multiplicatively independent here if
\[
Q^a\ne D^b\quad\text{for every }a,b\in\N_{>0}.
\]
This is a condition on their full prime-exponent vectors, not just their
gcd, their prime supports, or their numerical sizes.
\end{definition}

% A return would be an arithmetic resonance, independent of the offset.
\begin{lemma}\label{Lem:Injective}
Suppose \(|q|,|d|\ge2\) are multiplicatively independent and \(y_0\ne0\).
Then the sequence \((|y_n|)_{n\ge0}\) is injective.
\end{lemma}
\begin{proof}
Suppose \(i<j\) and \(|y_i|=|y_j|\). Apply Lemma~\ref{Lem:Product}
to the actual orbit segment starting at time \(i\), of length \(j-i\).
With its counts \(a,b\), absolute values and cancellation of
\(|y_i|>0\) give
\[
|d|^b=|q|^a,\qquad a+b=j-i>0.
\]
Neither exponent can vanish: if \(a=0\) then \(b>0\) and \(|d|^b>1\),
contradicting the equality; if \(b=0\) the same argument uses \(|q|^a>1\).
Both positive exponents therefore contradict independence.
\end{proof}

% Discreteness supplies full-sequence escape rather than an unbounded subsequence.
\begin{theorem}\label{Thm:Escape}
Let \(q,c,d\in\Z\), \(|q|,|d|\ge2\), and suppose \(|q|,|d|\) are
multiplicatively independent. For every nonreal \(z\in R\),
\[
\forall B\in\R\ \exists K\in\N\ \forall n\ge K:\quad |T^n(z)|>B.
\]
In particular every finite raw cycle lies on the real integer axis.
\end{theorem}
\begin{proof}
A nonreal seed has \(y_0\ne0\). For any integer \(H\ge0\), injectivity
allows at most \(H+1\) indices with \(|y_n|\le H\), because their values
belong to \(\{0,\ldots,H\}\). The indices in this finite set have a last
element if the set is nonempty. Taking \(K\) beyond that last element
(or \(K=0\) when empty) proves \(|y_n|\to\infty\) along the full sequence.
Now
\[
|T^n(z)|\ge|\operatorname{Im}T^n(z)|
=(\operatorname{Im}\theta)|y_n|,
\]
with a fixed strictly positive factor. For each real \(B\), choose a
nonnegative integer \(H\) with \((\operatorname{Im}\theta)H>B\), then
apply the preceding tail statement. A finite cycle would repeat a
transverse magnitude, so it cannot contain a nonreal point.
\end{proof}

\begin{remark}
The argument bounds the number of visits to each bounded transverse set,
not the last visiting time. It gives no uniform effective escape-time
bound in \(q,c,d,z,B\). Injectivity alone does not imply monotone growth.
The real axis is invariant and retains exactly its original signed
integer map; the theorem does not decide that remaining integer problem.
\end{remark}

\section{Finite arithmetic certificates and negative controls}

% A small determinant replaces a quantification over all positive exponents.
\begin{lemma}\label{Lem:Certificate}
Coprime \(Q,D\ge2\) are multiplicatively independent. More generally,
if two primes \(p,r\) satisfy
\[
v_p(Q)v_r(D)-v_r(Q)v_p(D)\ne0,
\]
then \(Q,D\) are multiplicatively independent.
\end{lemma}
\begin{proof}
If \(Q^a=D^b\), coprimality gives
\(\gcd(Q^a,D^b)=1\), whereas equality makes this gcd \(D^b>1\), a
contradiction. For the second criterion, take valuations at \(p\) and
\(r\): \(a v_p(Q)=b v_p(D)\) and
\(a v_r(Q)=b v_r(D)\). Cross multiplication and subtraction yield
\(a(v_p(Q)v_r(D)-v_r(Q)v_p(D))=0\). A positive \(a\) and a nonzero
minor make that impossible.
\end{proof}

For \(Q=12=2^2 3\), \(D=18=2\,3^2\), the exponent minor is
\(2\cdot2-1\cdot1=3\ne0\). Thus every nonreal seed escapes for
\((q,c,d)=(12,c,18)\), for every integer \(c\), even though the two
coefficients have identical prime support and gcd six.

% The escape mechanism must permit infinitely many legal divisions.
\begin{lemma}\label{Lem:InfiniteDivisions}
For \((q,c,d)=(12,0,18)\) and seed \(\theta\), every \(j\ge0\) has the
three-tick block
\[
8^j\theta\longmapsto12\,8^j\theta
\longmapsto144\,8^j\theta\longmapsto8^{j+1}\theta.
\]
In particular the orbit has infinitely many divisions, although it escapes.
\end{lemma}
\begin{proof}
The first transverse coefficient has no factor three, so eighteen does not
divide it. The second has exactly one factor three, so eighteen does not
divide it either. The third is divisible by eighteen, and its quotient
is \(8^{j+1}\). The real coordinate is zero throughout. The blocks
therefore have exactly their stated guards, and induction gives
\(T^{3j}(\theta)=8^j\theta\).
\end{proof}

% Dependence can support genuinely nonreal recurrent raw dynamics.
\begin{lemma}\label{Lem:Dependent}
For \((q,c,d)=(4,0,8)\), the seed \(\theta\) has primitive raw period five:
\[
\theta\longmapsto4\theta\longmapsto16\theta
\longmapsto2\theta\longmapsto8\theta\longmapsto\theta.
\]
\end{lemma}
\begin{proof}
Eight does not divide transverse coefficients one or four, so the first
two ticks multiply by four. It divides sixteen, so the next tick gives
two; two is not divisible by eight, so the next gives eight; the last tick
divides eight. The five listed transverse coefficients are distinct.
Also \(4^3=8^2\), exactly the excluded positive-power resonance.
\end{proof}

\section{A dynamical conjugacy, not an identification of rings}

% Scalar dynamics forgets the multiplication table without forgetting its guards.
\begin{theorem}\label{Thm:Transfer}
Let \(R'=\Z[\theta']\) be another imaginary quadratic coordinate order.
The coefficient matching
\[
\Phi(x+y\theta)=x+y\theta'
\]
is a bijection and satisfies
\(\Phi\circ T_{q,c,d}=T'_{q,c,d}\circ\Phi\).
It retains the original raw clock and the complete return-time set.
It need not preserve ring multiplication.
\end{theorem}
\begin{proof}
The inverse matches coefficients in the opposite direction.
Lemma~\ref{Lem:Guard} gives the same two integer divisibility tests in
both orders and the same scalar affine and division formulas. Thus the
one-tick identity holds in both branches; induction gives the identity
at every raw time. Bijectivity reflects as well as preserves each return.
For the last assertion, match the Gaussian basis element \(i\) to the
Eisenstein basis element \(\omega\). Then
\(\Phi(i^2)=\Phi(-1)=-1\), whereas
\(\Phi(i)^2=\omega^2=-1-\omega\ne-1\).
\end{proof}

\section{Reduction record and verification scope}

\paragraph{Original obligation.}
Classify the off-axis raw behavior for independent scalar parameters on
every specified imaginary quadratic order, retaining arbitrary real
integral offset and the full raw clock.
\paragraph{Replacement obligations.}
Prove actual ring-to-coordinate recovery; derive the transverse prefix
identity; exclude repeated magnitude by arithmetic independence; use the
local finiteness of the nonnegative integers and recover the complex norm.
\paragraph{Recovery.}
Lemma~\ref{Lem:Guard} preserves each guard, quotient and raw state.
Theorem~\ref{Thm:Escape} quantifies every real norm bound and all later
raw times. No unproved last-division assertion enters recovery.
\paragraph{Difficulty decrease.}
The two-coordinate adaptive guard no longer has to be solved. Its only
needed effect is an exactly integral division of the transverse coordinate.
The residual obstruction is a coefficient resonance, refuted by a finite
prime-exponent minor on each certified parameter pair, followed by finite
counting at each fixed height. This is a solved parameter subclass, not a
reduction of the remaining axis or dependent-scalar classification.
\paragraph{Preservation conditions.}
Scalar integral \(q,c,d\); \(|q|,|d|\ge2\); independence of absolute
magnitudes; a free integral quadratic basis with positive imaginary part;
and a nonzero transverse coefficient. Nonreal offsets, nonscalar divisors
and residue-wise quotient maps are not included.

\paragraph{Formal correspondence.}
The complete proof is in
\texttt{QCD.\allowbreak Complex\allowbreak Lattice.\allowbreak 001/\allowbreak lean/\allowbreak Scalar\allowbreak Escape.\allowbreak lean}.
Lemma~\ref{Lem:Guard} corresponds to
\texttt{Scalar\allowbreak Divides}, \texttt{Scalar\allowbreak Quotient}, \texttt{Scalar\allowbreak Unique},
\texttt{Ring\allowbreak Coordinates} and \texttt{Ring\allowbreak Coordinates\allowbreak Iter}.
Lemma~\ref{Lem:Product} is \texttt{Transverse\allowbreak Product} and
\texttt{Transverse\allowbreak Nonzero}; Lemma~\ref{Lem:Injective} is
\texttt{No\allowbreak Magnitude\allowbreak Return} and \texttt{Magnitude\allowbreak Injective};
Theorem~\ref{Thm:Escape} is proved by \texttt{Off\allowbreak Axis\allowbreak Escape} and
\texttt{Complex\allowbreak Off\allowbreak Axis\allowbreak Escape}. These use the complex embedding and
norm bound checked in \texttt{Complex\allowbreak Lattice.\allowbreak lean}.
Lemma~\ref{Lem:Certificate} is \texttt{Coprime\allowbreak Independent} and
\texttt{Exponent\allowbreak Independent}; the shared-support certificate is
\texttt{Twelve\allowbreak Eighteen\allowbreak Independent}.
Lemmas~\ref{Lem:InfiniteDivisions} and \ref{Lem:Dependent} are
\texttt{Common\allowbreak Support\allowbreak Block}, \texttt{Common\allowbreak Support\allowbreak Iterate} and
\texttt{Dependent\allowbreak Five\allowbreak Cycle}.
The proof of Theorem~\ref{Thm:Transfer} uses \texttt{Order\allowbreak Transfer\allowbreak Inverse},
\texttt{Order\allowbreak Transfer\allowbreak Raw} and \texttt{Order\allowbreak Transfer\allowbreak Iterate}.
The failure of multiplicativity is checked by
\texttt{Order\allowbreak Transfer\allowbreak Not\allowbreak Multiplicative}.

\paragraph{Computational scope and prior conventions.}
The Go survey produces additional scalar-family pictures and verifies
bounded signed arithmetic against independent exact references. Its
prime-exponent certificate accepts magnitudes at most \(65536\); this
implementation range is not a hypothesis of the Lean theorem.
A bounded trajectory cutoff is not itself an escape proof.
The residue-wise Gaussian convention described by Kucinski
\cite{Kucinski00} instead uses residue-dependent affine numerators followed
by division. Its nonreal cycles are not counterexamples to this raw map.
Only that definition and distinction are used here; no external global
classification theorem is imported.

\bibliographystyle{amsalpha}
\bibliography{ScalarEscape}

\end{document}
